\documentclass[a4paper,10pt]{article}

\usepackage[czech,english]{babel}
% Fonts %
\usepackage[T1]{fontenc}
\usepackage{FiraSans}
\renewcommand*\oldstylenums[1]{{\firaoldstyle #1}}
\usepackage{PTSerif}

% Colors %
\usepackage[dvipsnames]{color}
\usepackage{xcolor}

% Page Layout %
\usepackage[margin=0.5in]{geometry}

% Fancy Headers %
\usepackage{fancyhdr}
\fancyhf{}
\cfoot{}
\rhead{}
\renewcommand{\headrulewidth}{0pt}
\setlength{\headheight}{16pt}

% Math
\usepackage{mathtools}
\usepackage{amssymb}
\usepackage{faktor}
\usepackage{import}
\usepackage{caption}
\usepackage{subcaption}
\usepackage{wrapfig}
\usepackage{enumitem}
\usepackage{tikz}
\usetikzlibrary{cd,positioning,babel,shapes}
\usepackage{tkz-base}
\usepackage{tkz-euclide}

% Theorems
\usepackage{thmtools}
\usepackage[thmmarks, amsmath, thref]{ntheorem}

% Table of Contents %
\usepackage{hyperref}
\hypersetup{
 colorlinks=true,
 linktoc=all,
 linkcolor=blue
}

% Tables %
\usepackage{booktabs}
\usepackage{tabularx}
\usepackage{array}
\newcolumntype{?}{!{\vrule width 1pt\hspace*{.7pt}\vrule width 1pt}}

% Patch for hyphens
\usepackage{regexpatch}
\makeatletter
% Change the `-` delimiter to an active character
\xpatchparametertext\@@@cmidrule{-}{\cA-}{}{}
\xpatchparametertext\@cline{-}{\cA-}{}{}
\makeatother

\newcolumntype{s}{>{\centering\arraybackslash}p{.4\textwidth}}

% Operators %
\DeclareMathOperator{\Ker}{Ker}
\DeclareMathOperator{\Img}{Im}
\DeclareMathOperator{\End}{End}
\DeclareMathOperator{\Aut}{Aut}
\DeclareMathOperator{\Inn}{Inn}

% Common operators %
\newcommand{\R}{\mathbb{R}}
\newcommand{\N}{\mathbb{N}}
\newcommand{\Z}{\mathbb{Z}}
\newcommand{\Q}{\mathbb{Q}}
\newcommand{\C}{\mathbb{C}}

\newcommand{\tr}{\textcolor{red}}
\newcommand{\tb}{\textcolor{blue}}
\newcommand{\tg}{\textcolor{green}}
\newcommand{\tm}{\textcolor{magenta}}
\newcommand{\tv}{\textcolor{violet}}
\newcommand{\tmar}{\textcolor[HTML]{e23636}}
\newcommand{\tav}{\textcolor[HTML]{2e67a0}}

\definecolor{marvel}{HTML}{e23636}
\definecolor{avatar}{HTML}{2e67a0}

% American Paragraph Skip %
\setlength{\parindent}{0pt}
\setlength{\parskip}{1em}

\renewcommand{\baselinestretch}{1.2}
\renewcommand{\arraystretch}{1.2}
% Document %
\pagestyle{fancy}

\newcolumntype{C}[1]{%
 >{\vbox to 5ex\bgroup\vfill\centering}%
 p{#1}%
 <{\egroup}}

\begin{document}

\thispagestyle{fancy}
\begin{minipage}{.3\textwidth}
 \includegraphics[width=\textwidth]{logo}
\end{minipage}
\hfill
\begin{minipage}{.69\textwidth}
 \centering
 \Large{\sffamily
  GYMNÁZIUM EVOLUTION JIŽNÍ MĚSTO\\
  Posudek vedoucího ročníkové práce\\
  2022/23
 }
\end{minipage}

\begin{center}
 \large{
 \begin{tabular}{ll}
  \textsf{Název práce:} &
  \textsf{
   Psychická příprava sportovce a její vliv na výkon
  }\\
  \textsf{Autor práce, třída:} &
  \textsf{
   Matyáš Bogač, 3.C
  }\\
  \textsf{Vedoucí práce:} &
  \textsf{
   Pavel Kejha
  }
 \end{tabular}
 }
\end{center}

\section*{\sffamily \centering Hodnocení ročníkové práce}

\begin{center}
 \begin{tabular}{c|c|c|c|c|c|c|c|c|c}
  \sffamily
  Úv & LR & Me & An & Zh & Jaz & Cit & For & LS & Celkem\\
  \toprule
  \normalfont
  4
  &
  4
  &
  2
  &
  3
  &
  4
  &
  2
  &
  2
  &
  2
  &
  3
  &
  26
\end{tabular}
\end{center}

\begin{center}
  \begin{tabular}{l|l?l|l?l|l}
   Úv & Úvod & An & Analýza & Cit & Citace\\
   LR & Literární rešerše & Zh & Zhodnocení & For & Formátování\\
   Me & Metoda & Jaz & Jazyk & LS & Logická struktura\\
  \end{tabular}
 \end{center}

\section*{\sffamily \centering Slovní hodnocení výsledné práce (klady a
nedostatky)}
Práce pojednává o psychické přípravě a jejího vlivu na výkon dvou sportovců odlišného sportovního odvětví (kolektivní sport a individuální sport). Teoretická část popisuje  psychologický trénink,  roli rodiče a techniky pro trénink psychiky. V praktické části se zaměřuje na porovnání psychické přípravy a jejího vlivu na výkon a výsledky u dvou sportovců - zápasníka MMA a brankáře ledního hokeje. Práce je vyvážená a dobře strukturovaná. 
Po gramatické stránce by se dalo práci vytknout pár drobností - stylistické chyby (čárka, překlepy apod.). 
Škoda jen, že výzkum trval tak krátkou dobu, za kterou se až tolik toho neukázalo co se týká psychologického tréninku a jeho vlivu na výkon. V práci by mohl být zahrnutý dotazník, který by mapoval zda sportovci využívají mentálních koučů a proč.


\section*{\sffamily \centering Doporučené otázky k obhajobě}

\begin{enumerate}
  \item Pracoval jsi někdy ty sám s mentálním koučem?
\item Od kolika let bys doporučil začít pracovat s mentálním koučem?
\item Jak často bys doporučil konzultaci s mentálním kaučem?
\end{enumerate}

\vfill

\textbf{Datum:
01.04.2023
}
\hfill
\textit{
  Pavel Kejha
}

\end{document}
